\documentclass[12pt]{article}
\usepackage[utf8]{inputenc}
\usepackage{amsmath,amsfonts,amssymb}
\usepackage[amsthm]{ntheorem} %amsthm but also no newline in lists!
\usepackage[a4paper,margin=2.5cm]{geometry} %greater control over layout
\usepackage{tikz}
\usepackage{pgf, pgfplots} %allows tikzpictures
\pgfplotsset{compat=1.15}
\usepackage{mathrsfs}
\usetikzlibrary{arrows}
\usetikzlibrary{matrix}
\usepackage{url} %urls easy
\usepackage[shortlabels]{enumitem} %allows greater customisability of enum
\usepackage{textcomp} %gets rid of not defining \perthousand and \micro somehow
\usepackage{gensymb} %symbol package
\usepackage{sectsty}
\usepackage{tabto} %absolute/relative horizontal jumps in text
\usepackage{perpage} %footnote counter resets each page
\usepackage[symbol]{footmisc} %use symbols instead of numbers for footnotes
\usepackage{float} %allows "H" parameter for images
\usepackage[htt]{hyphenat} %tt text breaks properly
\usepackage{soul} %strikethrough text available
\usepackage[savepos]{zref} %allows absolute positioning on page
\usepackage{calc} %allows easy manipulation of lengths arithmetic
\usepackage{hyperref} %for hyperlinks in text. Package should be loaded last.

\hypersetup{
    colorlinks,
    linkcolor={red!50!black},
    citecolor={blue!50!black},
    urlcolor={blue!80!black}
}

\newcommand*{\currpos}[1]{%
  \zsavepos{#1}%
  %(\zposx{#1}sp, \zposy{#1}sp) =
  %(\the\dimexpr\zposx{#1}sp\relax, \the\dimexpr\zposy{#1}sp\relax)%
}

\begin{document}
	\pagenumbering{gobble}
		\begin{flushleft}
		\includegraphics[scale=0.7]{../../MODS_text_right.eps}\\
		\end{flushleft}
		\vspace{15pt}
		\noindent \LARGE\textbf{2020 August Intermediate Contest}\\[10pt]
		\noindent \normalsize{Saturday and Sunday, 8--9 August 2020}\\
		\noindent \makebox[\linewidth]{\rule{\textwidth}{0.4pt}}
			
	\normalsize
	
	\paragraph{\currpos{one}Problem 1.} For a positive integer $n$, let $m_n$ be the number formed by writing the digits of $n$ in ascending order, and $M_n$ be the number formed by writing the digits in descending order. For example, $m_{121} = 112, m_{420} = 24, M_{1337} = 7331$.\\[7pt]	
	\noindent Do there exist infinitely many positive integers $n$ such that $M_n - m_n = n$?
	
	\paragraph{\currpos{two}Problem 2.} Let \(\mathbb{R}^+\) denote the set of positive real numbers. Find all functions $f: \mathbb{R}^+ \to \mathbb{R}^+$ such that for all $x, y, z \in \mathbb{R}^+$ satisfying $x + y + z = xyz,$ $$f(xy) + f(xz) + f(yz) = 3.$$
	
	\paragraph{\currpos{three}Problem 3.} Let \(S_0\) be a set of 2020 real numbers. The sequence of sets \(S_0, S_1, \dots, S_n\) is called \emph{culled} if for all \(i \in \{0, 1, \ldots, n-1\}\),
	\begin{itemize}
		\item \(S_{i+1} \subset S_i\)
		\item No element \(x \in S_i\) with \(x \notin S_{i+1}\) is closer to the average of all elements in \(S_i\) than any element in \(S_{i+1}\).
	\end{itemize}
	An element \(b \in S_0\) is called \emph{basic} if there exists a culled sequence of sets such that \(S_n = \{b\}\). What is the least possible number of basic elements in \(S_0\)?
	
	\paragraph{\currpos{four}Problem 4.} Let \(\Gamma_1\) and \(\Gamma_2\) be two disjoint circles whose interiors do not intersect. Points \(X_1\) and \(X_2\) are on \(\Gamma_1\) and \(\Gamma_2\) respectively such that \(X_1X_2\) is an internal common tangent to \(\Gamma_1\) and \(\Gamma_2\). The circle $\Omega$ with diameter $X_1X_2$ intersects $\Gamma_1$ and $\Gamma_2$ again at $Y_1 \neq X_1$ and $Y_2 \neq X_2$ respectively. Suppose that $X_1Y_1$ and $X_2Y_2$ intersect at $P$. Prove that the perpendiculars from $P$ to the common external tangents of $\Gamma_1$ and $\Gamma_2$ are tangent to $\Omega$.
	
	\noindent \currpos{five} %\\[10pt]
	
	\vfill
	
	\noindent \textit{Language: English \hfill Time: 4 hours}\\
	\null \hfill \textit{Each problem is worth 7 points}\\
	
	\noindent \makebox[\linewidth]{\rule{\textwidth}{0.4pt}}	
	
	\noindent \footnotesize{\textit{Discord Server Invite Link: }\texttt{\href{https://discord.gg/94UnnAG}{discord.gg/94UnnAG}} \hfill \textit{Please support us on \href{https://bit.ly/MODS_patreon}{Patreon}!}}
	
	\begin{tikzpicture}[remember picture, overlay, shift=(current page.south west)]
    %top strip
    \draw[draw=none, fill=black, opacity=0.0] (current page.north west) rectangle (\paperwidth, \zposy{one}sp+18pt);
    %bottom strip
    \draw[draw=none, fill=black, opacity=0.0] (0,0) rectangle (\paperwidth, \zposy{five}sp+18pt);
	%left strip
    \draw[draw=none, fill=black, opacity=0.0] (0,\zposy{one}sp+18pt) rectangle (\zposx{one}sp-10pt,\zposy{five}sp+18pt );
    %right strip
    \draw[draw=none, fill=black, opacity=0.0] (\paperwidth-\zposx{one}sp+10pt,\zposy{five}sp+18pt) rectangle (\paperwidth, \zposy{one}sp+18pt);
    %P1
    \draw[draw=none, fill=black, opacity=0.0] (\zposx{one}sp-10pt,\zposy{one}sp+18pt) rectangle (\paperwidth-\zposx{one}sp+10pt, \zposy{two}sp+18pt);
    %P2
    \draw[draw=none, fill=black, opacity=0.0] (\zposx{two}sp-10pt,\zposy{two}sp+18pt) rectangle (\paperwidth-\zposx{two}sp+10pt, \zposy{three}sp+18pt);
    %P3
    \draw[draw=none, fill=black, opacity=0.0] (\zposx{three}sp-10pt,\zposy{three}sp+18pt) rectangle (\paperwidth-\zposx{three}sp+10pt, \zposy{four}sp+18pt);
    %P4
    \draw[draw=none, fill=black, opacity=0.0] (\zposx{four}sp-10pt,\zposy{four}sp+18pt) rectangle (\paperwidth-\zposx{four}sp+10pt, \zposy{five}sp+18pt);
	\end{tikzpicture}
		
		
		
	
	
\end{document}